\section{Observabilidade sem sujar o domínio}
\label{sec:observabilidade-quarkus}

Observabilidade é outro detalhe técnico: a regra de fechar pedido não
muda porque alguém quer um gráfico. As três peças abaixo entram por
extensão e ficam nos adaptadores e em \texttt{config}.

\paragraph{Health checks.}
O guia \citetitle{doc:quarkus-health} expõe \texttt{/q/health/live}
(\enquote{The application is up and running}),
\texttt{/q/health/ready} (pronta para atender), \texttt{/q/health/started}
e \texttt{/q/health}, que agrega todos. Um check é um bean que implementa
\texttt{HealthCheck} anotado com \texttt{@Liveness} ou \texttt{@Readiness}.
Extensões registram checks sozinhas: o guia cita a
\texttt{quarkus-agroal}, que registra um readiness check para cada
datasource. Coloque checks próprios em \texttt{adaptador} (por exemplo,
\texttt{adaptador.saida.cobranca.CobrancaHealth}), perto do recurso que
eles verificam.

Liveness e readiness respondem a perguntas diferentes, e confundi-las
custa caro num orquestrador: um liveness que depende do banco faz o
container ser reiniciado quando o banco cai, o que não conserta o
banco. Use dependências externas no readiness.

\paragraph{Métricas.}
O guia \citetitle{doc:quarkus-micrometer} diz que \enquote{Micrometer is
the recommended approach to metrics for Quarkus}; com a extensão
\texttt{micrometer-registry-prometheus}, as métricas saem em
\texttt{/q/metrics}. As anotações \texttt{@Timed} e \texttt{@Counted}
funcionam em métodos de beans. Ponha-as no resource ou no adaptador;
se uma métrica for de negócio (pedidos fechados por hora), publique um
evento de domínio e conte-o num adaptador, em vez de importar Micrometer
em \texttt{dominio}.

\paragraph{Traces.}
O guia \citetitle{doc:quarkus-opentelemetry} diz que, com a dependência
adicionada, a extensão gera dados de tracing por padrão e os exporta em
lote por OTLP/gRPC para \texttt{http://localhost:4317}. Métricas e logs
via OpenTelemetry são \emph{tech preview} e vêm desligados; ligam-se em
build time com \texttt{quarkus.otel.metrics.enabled=true} e
\texttt{quarkus.otel.logs.enabled=true}.

\paragraph{E a JVM?}
Essas três peças mostram a aplicação. Para ver a JVM (pausas de GC,
alocação, pinning de virtual threads, compilação), as ferramentas são as
da seção~\ref{sec:ferramentas}: JFR, \texttt{jcmd} e JMC. Os dois níveis
se completam: o trace mostra \emph{qual} requisição ficou lenta; o JFR
mostra \emph{por quê}, por exemplo uma pausa de GC no mesmo instante.
